\section*{Chapitre \ref{ch:b2:geological-hazards} --- Aléas géologiques et risque}

\begin{solution}{exo:b2:geological-hazards:1}
(a) $1 - e^{-30/100} = 0.259$. (b) $e^{-100/100} = e^{-1} = 0.368$ : plus d'un
siècle sur trois s'écoule sans une telle crue. (c) $T = -50/\ln 0.95 = 975$ ans.
\end{solution}

\begin{solution}{exo:b2:geological-hazards:2}
(a) Ville A~: $(1/200) \times 5000 \times 0.02 = 0.5$ bâtiment par an~; ville B~: $(1/50) \times 800
\times 0.10 = 1.6$. (b) Ville B, avec 6.25 fois moins de bâtiments exposés. (c) Son aléa est 4 fois plus fréquent et
sa vulnérabilité 5 fois plus élevée~: ensemble un facteur 20, qui l'emporte sur la
moindre \omterm{def:b2:geological-hazards:risk}{exposition}.
\end{solution}

\begin{solution}{exo:b2:geological-hazards:3}
(a) $\sqrt{9.81 \times 4000} = \qty{198}{m.s^{-1}}$~; $\sqrt{9.81 \times 200} = \qty{44.3}{m.s^{-1}}$~;
$\sqrt{9.81 \times 10} = \qty{9.90}{m.s^{-1}}$. (b) $\num{1.5e6}/198.1 = \qty{7570}{s}$, soit
\qty{2.10}{h}. (c) $800/3.6 = \qty{222.2}{m.s^{-1}}$, et $h = c^2/g = 222.2^2/9.81 = \qty{5030}{m}$.
\end{solution}

\begin{solution}{exo:b2:geological-hazards:4}
(a) $\tan\ang{33}/\tan\ang{25} = 0.6494/0.4663 = 1.39$. (b) $F = 1$ lorsque $\beta = \varphi' =
\ang{33}$. (c) Avec $c' = 0$, $F = (\cos^2\beta - r_u)\tan\varphi'/(\sin\beta\cos\beta)$, et $\cos^2\beta
- r_u = \cos^2\beta\,(1 - \gamma_w/\gamma)$, donc $F = (1 - \gamma_w/\gamma)\tan\varphi'/\tan\beta$~:
l'eau réduit de moitié le frottement. $(1 - 9.81/20) \times 1.393 = 0.510 \times 1.393 = 0.710$~: la
pente rompt.
\end{solution}

\begin{solution}{exo:b2:geological-hazards:5}
(a) $696/50 = 13.9$ par an. (b) $e^{-30/26.2} = 0.318$, contre 0.304 observé~:
proche. La valeur observée est un peu plus basse parce que les répliques et les
doublets regroupés ajoutent des intervalles courts. (c) $1 - e^{-7/26.2} = 0.234$.
\end{solution}

\begin{solution}{exo:b2:geological-hazards:6}
(a) Le terme de magnitude vaut $-0.105 + 0.229 \times 0.5 = 0.0095$. Pour $d = 0$, 10, 30 et
\qty{100}{km}, $r = 5.57$, 11.45, 30.51 et \qty{100.2}{km}, $\log_{10}\mathrm{PGA} = -0.571$,
$-0.814$, $-1.145$ et $-1.547$, et la PGA médiane vaut 0.269, 0.153, 0.0715 et \qty{0.0284}{g}.
(b) $0.153 \times 10^{0.251} = 0.153 \times 1.78 = \qty{0.272}{g}$. (c) Un écart type de
part et d'autre~: $0.153/1.70 = \qty{0.0900}{g}$ à $0.153 \times 1.70 = \qty{0.260}{g}$.
\end{solution}

\begin{solution}{exo:b2:geological-hazards:7}
(a) Classés~: 420 ($T = 13/1 = 13$ ans), 305 (6.5), 268 (4.33), 251 (3.25). (b) $305 + (420 -
305)\ln(10/6.5)/\ln(13/6.5) = 305 + 115 \times 0.621 = \qty{376}{m^3.s^{-1}}$. (c) Le plus grand pic
est placé à 13 ans~; toute valeur à 100 ans extrapole loin au-delà de l'enregistrement, avec une loi des
extrêmes dont le choix et les paramètres sont à peine contraints par les 12 points.
\end{solution}

\begin{solution}{exo:b2:geological-hazards:8}
(a) $0.6 \times 0.4 \times 0.5 \times 0.3 = 0.036$. (b) $P(\text{éruption}) = 0.6 \times 0.4 = 0.24$~;
sans coulée dans la vallée, $0.24 - 0.036 = 0.204$. (c) $0.6 \times 0.7 \times 0.5 \times 0.3 =
0.063$.
\end{solution}

\begin{solution}{exo:b2:geological-hazards:9}
Ici $\sin\beta\cos\beta = 0.4698$, $\cos^2\beta = 0.6710$ et $\tan\ang{30}/\tan\ang{35} =
0.8245$. (a) $F = c'/(\gamma z\sin\beta\cos\beta) + 0.8245 = 1$ donne $z_c = 8/(19 \times 0.4698
\times 0.1755) = \qty{5.11}{m}$. (b) $8/(19 \times 2 \times 0.4698) + 0.8245 = 0.448 + 0.825 =
1.27$. (c) $r_u = 9.81 \times 0.6710/19 = 0.3465$~; $F = [8 + 38 \times (0.6710 - 0.3465) \times
0.5774]/(38 \times 0.4698) = 15.12/17.85 = 0.847$~: la couche saturée rompt.
\end{solution}

\begin{solution}{exo:b2:geological-hazards:10}
(a) $\Pr(\text{aucun dans } [0, s+t] \mid \text{aucun dans } [0, s]) = e^{-(s+t)/T}/e^{-s/T} = e^{-t/T}$.
(b) Les 400 années calmes ne changent rien~: $1 - e^{-50/475} = 0.0999$, soit \qty{10}{\percent}. (c) Cela
contredit l'hypothèse selon laquelle le taux ne dépend pas du temps écoulé~: après le rebond, la contrainte
doit se reconstituer, si bien que le taux est faible juste après un séisme et croît ensuite. Une zone source
regroupe de nombreuses failles à différents stades de leur cycle~; la somme de nombreuses suites indépendantes
d'événements se comporte presque comme un processus de Poisson, et le taux est tout ce qu'on connaît de chacune.
\end{solution}

\begin{solution}{exo:b2:geological-hazards:11}
(a) $A = 0.5 \times (4000/10)^{1/4} = 0.5 \times 4.47 = \qty{2.24}{m}$. (b) $\lambda = cT$~:
$198.1 \times 1200 = \qty{238}{km}$ en eau profonde, $9.90 \times 1200 = \qty{11.9}{km}$
dans \qty{10}{m} d'eau. (c) L'amplitude n'est plus petite devant la profondeur
(\qty{2.2}{m} dans \qty{10}{m}), si bien que l'onde se raidit et peut déferler~; une
partie de l'énergie est réfléchie ou concentrée par la forme du fond marin et des
baies~; sur terre, la \omterm{def:b2:geological-hazards:tsunami}{hauteur de run-up} dépend de la pente et de la rugosité du terrain.
\end{solution}

\begin{solution}{exo:b2:geological-hazards:12}
(a) $10/5.8 + 4 = 1.72 + 4 = \qty{5.72}{s}$. (b) L'onde S atteint \qty{50}{km} après $50/3.2 =
\qty{15.6}{s}$~: $15.6 - 5.72 = \qty{9.90}{s}$ d'alerte~; à \qty{100}{km}, $31.3 - 5.72 =
\qty{25.5}{s}$. (c) $3.2 \times 5.72 = \qty{18.3}{km}$.
\end{solution}

\begin{solution}{pb:b2:geological-hazards:1}
Les logarithmes sont gardés à quatre décimales dans les étapes intermédiaires.
\textbf{1.} $10^{4-5} = 0.1$, $10^{4-6} = 0.01$ et $10^{4-7} = 0.001$ par an.
\textbf{2.} 10, 100 et 1000 ans.
\textbf{3.} $10^{4-7.5} = \num{3.16e-4}$ par an~; $0.001 - \num{3.16e-4} = \num{6.84e-4}$ par an.
\textbf{4.} $0.1 - 0.01 = 0.09$ et $0.01 - 0.001 = 0.009$ par an.
\textbf{5.} $1 - e^{-0.01 \times 50} = 0.393$.
\textbf{6.} $1 - e^{-0.001 \times 50} = 0.0488$.
\textbf{7.} $r = \sqrt{10^2 + 5.57^2} = \qty{11.45}{km}$~; $\log_{10} r = 1.0588$ et $-0.778 \times
1.0588 = -0.8237$.
\textbf{8.} $\mu = -0.105 + 0.229(M - 6) - 0.8237$~: $-1.0432$,
$-0.8142$ et $-0.6424$ ; médianes 0.0905, 0.153 et \qty{0.228}{g}.
\textbf{9.} $10^{-0.8142 + 0.230} = \qty{0.260}{g}$.
\textbf{10.} $z = (\log_{10} 0.3 + 0.8142)/0.230 = (-0.5229 + 0.8142)/0.230 = 1.266$, et $1 -
\Phi(1.266) = 0.103$.
\textbf{11.} $-0.105 + 0.229 \times 1.5 - 0.8237 = -0.5852$~: \qty{0.260}{g}.
\textbf{12.} Les probabilités de dépasser \qty{0.2}{g} ($\log_{10} 0.2 = -0.6990$) valent 0.0672
($z = 1.497$), 0.308 ($z = 0.501$) et 0.597 ($z = -0.246$)~: $\lambda = 0.09 \times 0.0672 + 0.009
\times 0.308 + \num{6.84e-4} \times 0.597 = \num{6.05e-3} + \num{2.77e-3} + \num{4.08e-4} =
\num{9.23e-3}$ par an.
\textbf{13.} Pour \qty{0.3}{g}~: $z = 2.262$, 1.266 et 0.520, probabilités 0.0118, 0.103 et
0.302~; $\lambda = \num{1.07e-3} + \num{9.24e-4} + \num{2.06e-4} = \num{2.20e-3}$ par an.
\textbf{14.} Pour \qty{0.4}{g} ($\log_{10} 0.4 = -0.3979$)~: $z = 2.805$, 1.810 et 1.063,
probabilités 0.00251, 0.0352 et 0.144~; $\lambda = \num{2.26e-4} + \num{3.17e-4} + \num{9.84e-5} =
\num{6.41e-4}$ par an.
\textbf{15.} $-\ln 0.9/50 = \num{2.107e-3}$ par an, soit une \omterm{def:b2:geological-hazards:return}{période de retour} de 475 ans.
\textbf{16.} Entre \qty{0.3}{g} (\num{2.20e-3}) et \qty{0.4}{g} (\num{6.41e-4}).
\textbf{17.} La pente vaut $\ln(6.41/22.0)/\ln(4/3) = -1.233/0.2877 = -4.29$, et $a = 0.3 \times
(2.107/2.20)^{-1/4.29} = 0.3 \times 1.0101 = \qty{0.303}{g}$.
\textbf{18.} À \qty{0.3}{g}, la classe 5--6 fournit \qty{48.5}{\percent} du taux, la classe
6--7 \qty{42.1}{\percent}, et 7--7.5 seulement \qty{9.4}{\percent}. Même le plus fort séisme
a une médiane de \qty{0.260}{g}, inférieure à \qty{0.3}{g}~: ce niveau n'est atteint que dans
la queue haute de la dispersion, et les séismes modérés fréquents en fournissent l'essentiel.
\textbf{19.} $4/(19 \times 3 \times 0.4330) + \tan\ang{34}/\tan\ang{30} = 0.162 + 1.168 = 1.33$.
\textbf{20.} $r_u = 9.81 \times 0.75/19 = 0.387$~; $F = [4 + 57 \times (0.75 - 0.387) \times
0.6745]/(57 \times 0.4330) = 17.95/24.68 = 0.727$~: la pente saturée rompt.
\textbf{21.} Avec la nappe à une fraction $m$ de l'épaisseur, $r_u = 0.387\,m$.
$F = 1$ lorsque $0.75 - 0.387\,m = (24.68 - 4)/(57 \times 0.6745) = 0.538$~:
$m = 0.548$, soit une nappe à \qty{55}{\percent} de la couche.
\textbf{22.} En divisant $F = 1$ par $\gamma z$~: $0.0702 + (0.75 - 0.4330\,k) \times 0.6745 = 0.4330 +
0.75\,k$, d'où $k_c = (0.0702 + 0.5059 - 0.4330)/(0.2921 + 0.75) = 0.1431/1.042 = 0.137$.
\textbf{23.} $\log_{10} 0.137 = -0.8633 = -0.105 + 0.229(M - 6) - 0.8237$ donne $M = 6.29$. De tels
séismes surviennent à $10^{4 - 6.29} - 10^{-3.5} = \num{4.81e-3}$ par an, soit environ une fois en
208 ans, et la PGA de 475 ans de la question~17 est plus de deux fois le \qty{0.137}{g} critique :
selon cette analyse, la pente sèche rompt à une secousse bien inférieure au mouvement de 475 ans.
\textbf{24.} La PGA ayant une probabilité de \qty{10}{\percent} d'être dépassée en
50 ans au niveau de la ville vaut \textbf{\qty[text-series-to-math=true]{0.303}{g}}.
\end{solution}
